\section{OpenAI Provider 与非 OpenAI Provider 的路径分叉}

\subsection{两条不同的代码路径}

Codex 源码中根据 Provider 类型走不同的 Compact 路径：

\begin{itemize}
    \item \textbf{OpenAI Provider}：调用专门的 \texttt{response.compact} API——这是一个专用的压缩接口，只需提交对话历史，即可获得压缩版内容。System Prompt 和 Handoff Prompt 在\textbf{远端}（OpenAI 服务器端）维护。
    \item \textbf{非 OpenAI Provider}：调用普通的 \texttt{response.create} 接口，在\textbf{本地}注入 Compact Prompt 和 Handoff Prompt。
\end{itemize}

\begin{knowledgebox}{response.compact API}
这是 OpenAI 提供的专用压缩接口。调用方只需提交对话历史（messages），API 返回压缩后的摘要。无需在客户端维护压缩提示词。
\end{knowledgebox}

\subsection{本章小结}

Provider 的不同决定了压缩提示词的维护位置：OpenAI 走远端 \texttt{response.compact}，其余走本地 \texttt{response.create} 并显式注入提示词。


